\section{Реализация, тестирование и демонстрация гибридной интеллектуальной системы}

\subsection{Выбор средств и технологий реализации}
% Обосновывается выбор языков программирования, сред разработки, библиотек,
% фреймворков, СУБД, средств работы с базой знаний, нейробионическими
% моделями, векторными хранилищами, API и средствами тестирования.
%
% Например, могут быть использованы:
% - Python, Prolog, Java/JVM;
% - RDFLib, Owlready2, OWL API, Apache Jena;
% - CLIPS, Drools, Experta;
% - PyTorch, TensorFlow/Keras, scikit-learn;
% - Hugging Face Transformers;
% - LangChain, LlamaIndex, Haystack, Semantic Kernel;
% - FAISS, Qdrant, Chroma, Weaviate;
% - FastAPI, Flask, Django REST Framework, Spring Boot;
% - pytest, средства интеграционного и end-to-end тестирования.
%
% Перечисляются только средства, действительно использованные в проекте,
% и обосновывается их соответствие спроектированной архитектуре.

\subsection{Структура программной реализации}
% Описывается фактическая структура исходного кода и программного проекта:
% модули, пакеты, классы, сервисы, файлы конфигурации, хранилища данных,
% база знаний, модели, точки входа и средства развертывания.
%
% Показывается соответствие структуры реализации логической модели
% и архитектуре гибридной интеллектуальной системы.

\subsection{Реализация пользовательского интерфейса}
% Описывается реализация пользовательского или программного интерфейса,
% используемого для запуска гибридной интеллектуальной системы.
%
% Показывается, каким образом пользователь вводит запрос, получает
% результат, объяснение, уровень уверенности, предупреждения и сообщения
% об ошибках. Приводится связь интерфейса с реализованными сквозными
% сценариями.

\subsection{Тестирование интеллектуальной системы}
% Проверяется совместная работа компонентов в типичных, граничных
% и ошибочных сценариях.
%
% Для каждого теста указываются входные данные, используемые знания
% и модели, ожидаемый результат, фактический результат, оценка уверенности,
% объяснение и результат обработки возможных рассогласований.
%
% Проверяется сквозной путь от пользовательского запроса через интерфейс,
% решатель задач, базу знаний и нейробионический компонент до итогового
% решения.

\subsection{Оценка качества, эффективности и объяснимости системы}
% Анализируются корректность, полнота, устойчивость, время отклика,
% вычислительная нагрузка, качество результатов нейробионического
% компонента, влияние объёма базы знаний и качества данных на работу
% гибридной системы.
%
% При возможности выполняется сравнительный анализ:
% - варианта, использующего только логико-символьный подход;
% - варианта, использующего только нейробионический подход;
% - разработанной гибридной интеллектуальной системы.
%
% Оценивается, каким образом гибридный подход влияет на качество решения,
% объяснимость, устойчивость и производительность.

\subsection{Демонстрация работы гибридной интеллектуальной системы}
% Приводится последовательность демонстрации программного прототипа
% для нескольких осмысленных сквозных сценариев.
%
% Демонстрация должна показать:
% - ввод пользовательского или системного запроса;
% - использование базы знаний;
% - работу логико-символьного компонента;
% - работу нейробионического компонента;
% - интеграцию и согласование результатов;
% - итоговое решение, объяснение, уровень уверенности и обработку ошибок.

\subsection{Анализ ограничений и направления развития}
% Формулируются выявленные ограничения:
% - границы применимости предметной области;
% - полнота и актуальность базы знаний;
% - ограничения алгоритмов решения задач;
% - качество и устойчивость нейробионической модели;
% - ограничения объяснимости;
% - производительность и масштабируемость;
% - зависимость от внешних сервисов или данных.
%
% Указываются конкретные направления дальнейшего развития:
% расширение базы знаний, улучшение модели, добавление новых сценариев,
% развитие интерфейса, повышение объяснимости, оптимизация производительности,
% внедрение дополнительных механизмов контроля качества или масштабирование.


\subsection{Выводы}
% Кратко фиксируются результаты реализации, компонентного, интеграционного
% и сквозного тестирования, оценки качества и демонстрации работы
% гибридной интеллектуальной системы.
